% Generated draft: requires mathematical and source review.
\documentclass[11pt]{article}
\usepackage[margin=25mm]{geometry}
\usepackage{fontspec}
\setmainfont{Latin Modern Roman}
\usepackage{amsmath,amssymb,mathtools,unicode-math}
\setmathfont{Latin Modern Math}
\usepackage{enumitem,needspace,xurl,hyperref}
\setlength{\parindent}{0pt}
\setlength{\parskip}{5pt}
\setlength{\emergencystretch}{3em}
\newcommand{\N}{\mathbb{N}}
\newcommand{\Z}{\mathbb{Z}}
\newcommand{\Q}{\mathbb{Q}}
\newcommand{\R}{\mathbb{R}}
\newcommand{\C}{\mathbb{C}}
\newcommand{\F}{\mathbb{F}}
\newcommand{\A}{\mathbb{A}}
\newcommand{\PP}{\mathbb{P}}
\newcommand{\E}{\mathbb{E}}
\newcommand{\eps}{\varepsilon}
\newcommand{\dd}{\,\mathrm{d}}
\begin{document}
\section*{UnsolvedMath ID: 30005132}
\textbf{Source problem code:} catalogue entry
\subsection{Definitions and mathematical statement}
Let \(p\) be a prime and let \(X_1,X_2\) be algebraic varieties over the finite field \(\F_p\). Write \(X_i^{\mathrm{Weil}}\) for the Weil object associated with \(X_i\), and \(\operatorname{DTM}(Y)\) for the Weil-motivic Tate category attached to an admissible object \(Y\), using the source's intended conventions. The supplied record does not define these constructions, their coefficients, or the categorical framework for the relative tensor product; consequently these symbols cannot be given complete definitions from the supplied evidence alone.

The intended question is whether, for every such pair \(X_1,X_2\), exterior product induces an equivalence
\[
\boxtimes:\operatorname{DTM}(X_1^{\mathrm{Weil}})
\otimes_{\operatorname{DTM}(\F_p)}
\operatorname{DTM}(X_2^{\mathrm{Weil}})
\longrightarrow
\operatorname{DTM}\!\left(X_1^{\mathrm{Weil}}\times_{\F_p}X_2^{\mathrm{Weil}}\right).
\]
Here the tensor product is intended to be relative to the displayed base category \(\operatorname{DTM}(\F_p)\). The necessary module structures, the meaning of the fiber product of Weil objects, and the existence and precise construction of this functor require the source's missing conventions. The target is retained exactly as supplied; no identification with \((X_1\times_{\F_p}X_2)^{\mathrm{Weil}}\) is assumed.

As background, the supplied excerpt labels an analogous exterior-product equivalence for constructible \(\ell\)-adic sheaves as Theorem 2, attributed to reference [HRS20]. This is supplied background attribution, not an independently verified theorem here; its coefficient conventions and precise hypotheses are absent. It does not resolve the displayed motivic question.
\subsection{Short English statement}
For every pair of algebraic varieties over \(\F_p\), does exterior product identify the relative tensor product of their Weil-motivic Tate categories with the category on the displayed fiber product of Weil objects? The supplied record labels this question open, but essential categorical definitions and conventions remain unspecified.
\subsection{Sources}
\begin{itemize}
\item Local registry: \texttt{data/unsolvedmath/problems.json}, supplied record with permanent UnsolvedMath ID \(30005132\), source problem code \texttt{catalogue entry}, titled ``K\"unneth Equivalence for Weil Motivic Tate Categories.'' No individual proposer is supplied; the proposed year is recorded as 2022.
\item Supplied report citation: \emph{Algebraic K-Theory (2023)}, Oberwolfach Reports, pp.~1313--1376, DOI \texttt{10.4171/owr/2022/24}. \url{https://doi.org/10.4171/owr/2022/24}. The record describes this collection as arising from workshop open-problem discussions. Its 2023 title date and 2022 DOI/proposed-year information are retained without reconciliation.
\item Supplied literature reference: Hemo--Richarz--Scholbach, ``sheaf Kunneth,'' \url{https://arxiv.org/abs/2012.02853}. The extracted report statement attributes its sheaf comparison theorem to [HRS20]; the supplied record associates the sheaf comparison with this link.
\item Literature assessment dated exactly \texttt{2026-08-22}: supplied dated curation, not independently verified. It labels the problem open, states that the analogous K\"unneth equivalence for constructible \(\ell\)-adic sheaves is known, and describes the Weil-motivic Tate equivalence as ongoing and not established in the required generality. The listed sources are the Hemo--Richarz--Scholbach link and the Oberwolfach report link above. The supplied curation expressly requires independent verification of time-sensitive status and cited claims.
\item The registry labels the statement \texttt{corrected\_verified} and reports a direct check of the official report, with extraction artifacts removed and the mathematical target preserved. These are supplied verification claims; no independent source consultation or mathematical verification is asserted in this draft.
\end{itemize}
\end{document}
